
\documentclass{article}
\usepackage{colortbl}
\usepackage{makecell}
\usepackage{multirow}
\usepackage{supertabular}

\begin{document}

\newcounter{utterance}

\centering \large Interaction Transcript for game `wordle', experiment `medium\_frequency\_words\_no\_clue\_no\_critic', episode 15 with refineonly{-}Qwen3.5{-}9B.
\vspace{24pt}

{ \footnotesize  \setcounter{utterance}{1}
\setlength{\tabcolsep}{0pt}
\begin{supertabular}{c@{$\;$}|p{.15\linewidth}@{}p{.15\linewidth}p{.15\linewidth}p{.15\linewidth}p{.15\linewidth}p{.15\linewidth}}
    \# & \multicolumn{2}{c}{Player} && \multicolumn{2}{c}{Game Master} \\
    \hline

    \theutterance \stepcounter{utterance}  
    & & & \multicolumn{4}{p{0.6\linewidth}}{
        \cellcolor[rgb]{0.9,0.9,0.9}{
            \makecell[{{p{\linewidth}}}]{
                \texttt{\tiny{[P1$\langle$GM]}}
                \texttt{You are a language wizard who likes to guess words by using the given rules.} \\
\\ 
\texttt{Welcome to Wordle! You have six attempts to guess the target word, a valid English word of five lowercase letters (a{-}z). Please use the tags "explanation:" and "guess:" to provide a concise explanation for each guess.} \\
\\ 
\texttt{For instance, if your guess is "apple", your response should be} \\
\texttt{explanation: this is a common five{-}letter English word, and I am starting my guess with this word.} \\
\texttt{guess: apple} \\
\\ 
\texttt{After each guess, your answer will be validated, and you will receive feedback indicating which letters are correct (green), which letters are correct but in the wrong position (yellow), and which letters are incorrect (red). This feedback can be useful in determining which letters to include or exclude in your next guess.} \\
\\ 
\texttt{For example, the feedback for "apple" might be:} \\
\texttt{guess\_feedback: a<yellow> p<yellow> p<green> l<yellow> e<red>} \\
\\ 
\texttt{The explanation should contain details about how the guess\_feedback is used to arrive at a new guess.} \\
\\ 
\texttt{Do not generate any other text.} \\
\\ 
\texttt{Let's begin with your first guess.} \\
            }
        }
    }
     \\ \\

    \theutterance \stepcounter{utterance}  
    & & \multicolumn{4}{p{0.6\linewidth}}{
        \cellcolor[rgb]{0.95,0.95,0.95}{
            \makecell[{{p{\linewidth}}}]{
                \texttt{\tiny{[P1$\leftrightarrow$P]}}
                \texttt{{[}DRAFT{]}} \\
\texttt{explanation: I am starting with "crane" because it contains common vowels (a, e) and frequent consonants (r, n) that appear in many five{-}letter words, allowing me to quickly eliminate or confirm several letters.} \\
\texttt{guess: crane} \\
            }
        }
    }
    & \\ \\

    \theutterance \stepcounter{utterance}  
    & \multicolumn{4}{p{0.6\linewidth}}{
        \cellcolor[rgb]{0.561,0.737,0.561}{
            \makecell[{{p{\linewidth}}}]{
                \texttt{\tiny{[P1$\rangle$GM]}}
                \texttt{explanation: "crane" is an excellent starting word because it contains common vowels (a, e) and frequent consonants (c, r, n) that appear in many five{-}letter words, allowing me to quickly eliminate or confirm several letters.} \\
\texttt{guess: crane} \\
            }
        }
    }
    & & \\ \\

    \theutterance \stepcounter{utterance}  
    & & & \multicolumn{2}{p{0.3\linewidth}}{
        \cellcolor[rgb]{0.9,0.9,0.9}{
            \makecell[{{p{\linewidth}}}]{
                \texttt{\tiny{[GM$|$GM]}}
                \texttt{attempts = 1} \\
\texttt{target\_word = liner} \\
\texttt{guess = crane} \\
\texttt{guess\_feedback = c<red> r<yellow> a<red> n<yellow> e<yellow>} \\
            }
        }
    }
    & & \\ \\

    \theutterance \stepcounter{utterance}  
    & & & \multicolumn{4}{p{0.6\linewidth}}{
        \cellcolor[rgb]{0.9,0.9,0.9}{
            \makecell[{{p{\linewidth}}}]{
                \texttt{\tiny{[P1$\langle$GM]}}
                \texttt{guess\_feedback: c<red> r<yellow> a<red> n<yellow> e<yellow>} \\
\\ 
\texttt{Provide your response only in this format.} \\
\texttt{explanation: details} \\
\texttt{guess: word} \\
            }
        }
    }
     \\ \\

    \theutterance \stepcounter{utterance}  
    & & \multicolumn{4}{p{0.6\linewidth}}{
        \cellcolor[rgb]{0.95,0.95,0.95}{
            \makecell[{{p{\linewidth}}}]{
                \texttt{\tiny{[P1$\leftrightarrow$P]}}
                \texttt{{[}DRAFT{]}} \\
\texttt{explanation: The feedback indicates that 'r', 'n', and 'e' are present but misplaced, while 'c' and 'a' are absent. Since 'e' is in the word but not at the end, and 'r' and 'n' are also misplaced, I will try to place 'e' in the second position and keep 'r' and 'n' in different spots to test their correct locations.} \\
\texttt{guess: never} \\
            }
        }
    }
    & \\ \\

    \theutterance \stepcounter{utterance}  
    & \multicolumn{4}{p{0.6\linewidth}}{
        \cellcolor[rgb]{0.561,0.737,0.561}{
            \makecell[{{p{\linewidth}}}]{
                \texttt{\tiny{[P1$\rangle$GM]}}
                \texttt{explanation: The feedback indicates that 'r', 'n', and 'e' are present but misplaced, while 'c' and 'a' are absent. Since 'e' is in the word but not at the end, and 'r' and 'n' are also misplaced, I will try to place 'e' in the second position and keep 'r' and 'n' in different spots to test their correct locations.} \\
\texttt{guess: never} \\
            }
        }
    }
    & & \\ \\

    \theutterance \stepcounter{utterance}  
    & & & \multicolumn{2}{p{0.3\linewidth}}{
        \cellcolor[rgb]{0.9,0.9,0.9}{
            \makecell[{{p{\linewidth}}}]{
                \texttt{\tiny{[GM$|$GM]}}
                \texttt{attempts = 2} \\
\texttt{target\_word = liner} \\
\texttt{guess = never} \\
\texttt{guess\_feedback = n<yellow> e<red> v<red> e<green> r<green>} \\
            }
        }
    }
    & & \\ \\

    \theutterance \stepcounter{utterance}  
    & & & \multicolumn{4}{p{0.6\linewidth}}{
        \cellcolor[rgb]{0.9,0.9,0.9}{
            \makecell[{{p{\linewidth}}}]{
                \texttt{\tiny{[P1$\langle$GM]}}
                \texttt{guess\_feedback: n<yellow> e<red> v<red> e<green> r<green>} \\
\\ 
\texttt{Provide your response only in this format.} \\
\texttt{explanation: details} \\
\texttt{guess: word} \\
            }
        }
    }
     \\ \\

    \theutterance \stepcounter{utterance}  
    & & \multicolumn{4}{p{0.6\linewidth}}{
        \cellcolor[rgb]{0.95,0.95,0.95}{
            \makecell[{{p{\linewidth}}}]{
                \texttt{\tiny{[P1$\leftrightarrow$P]}}
                \texttt{{[}DRAFT{]}} \\
\texttt{explanation: The feedback shows 'n' is misplaced, 'e' is incorrect in the second position but correct in the fourth, and 'v' is absent. Since 'n', 'e', and 'r' are confirmed letters with 'e' fixed at position 4, I need to find a word with 'n' elsewhere, 'e' at the end, and 'r' in a new spot, while avoiding 'v'.} \\
\texttt{guess: never} \\
            }
        }
    }
    & \\ \\

    \theutterance \stepcounter{utterance}  
    & \multicolumn{4}{p{0.6\linewidth}}{
        \cellcolor[rgb]{0.561,0.737,0.561}{
            \makecell[{{p{\linewidth}}}]{
                \texttt{\tiny{[P1$\rangle$GM]}}
                \texttt{explanation: The previous feedback for "never" showed 'n' (yellow), 'e' (red at pos 2, green at pos 4), and 'v' (red). The planned guess "never" is identical to the previous one, which is invalid as it provides no new information. Since 'e' is confirmed at position 4, 'n' is present but not at position 1, and 'r' is present (from the first turn's "crane" feedback where 'r' was yellow, and confirmed here as part of the word structure if we assume standard Wordle logic where repeated letters in feedback imply presence, though strictly 'r' wasn't in "never", wait: in "crane", 'r' was yellow. In "never", 'r' is at the end. The feedback for "never" was n<yellow> e<red> v<red> e<green> r<green>. This confirms 'n', 'e', 'r' are in the word. 'e' is at pos 4. 'n' is not at pos 1. 'r' is at pos 5 (green). So the word ends in 'er'. 'n' is somewhere else. 'v' is out. I need a new word ending in 'er' with 'n' not at the start.} \\
\texttt{guess: inner} \\
            }
        }
    }
    & & \\ \\

    \theutterance \stepcounter{utterance}  
    & & & \multicolumn{2}{p{0.3\linewidth}}{
        \cellcolor[rgb]{0.9,0.9,0.9}{
            \makecell[{{p{\linewidth}}}]{
                \texttt{\tiny{[GM$|$GM]}}
                \texttt{attempts = 3} \\
\texttt{target\_word = liner} \\
\texttt{guess = inner} \\
\texttt{guess\_feedback = i<yellow> n<red> n<green> e<green> r<green>} \\
            }
        }
    }
    & & \\ \\

    \theutterance \stepcounter{utterance}  
    & & & \multicolumn{4}{p{0.6\linewidth}}{
        \cellcolor[rgb]{0.9,0.9,0.9}{
            \makecell[{{p{\linewidth}}}]{
                \texttt{\tiny{[P1$\langle$GM]}}
                \texttt{guess\_feedback: i<yellow> n<red> n<green> e<green> r<green>} \\
\\ 
\texttt{Provide your response only in this format.} \\
\texttt{explanation: details} \\
\texttt{guess: word} \\
            }
        }
    }
     \\ \\

    \theutterance \stepcounter{utterance}  
    & & \multicolumn{4}{p{0.6\linewidth}}{
        \cellcolor[rgb]{0.95,0.95,0.95}{
            \makecell[{{p{\linewidth}}}]{
                \texttt{\tiny{[P1$\leftrightarrow$P]}}
                \texttt{{[}DRAFT{]}} \\
\texttt{explanation: The feedback for "inner" confirms the word is "inner" itself: 'i' is yellow (pos 1 is wrong, but 'i' is in the word), 'n' at pos 2 is red (wrong spot), 'n' at pos 3 is green (correct), 'e' at pos 4 is green (correct), and 'r' at pos 5 is green (correct). Since all letters match the target word's positions except for the first letter 'i' which is actually correct in the word "inner" but marked yellow? Wait, if the word is "inner", then 'i' at pos 1 should be green. The feedback says 'i'<yellow>. This implies the target word is NOT "inner". Let's re{-}evaluate. Target has 'n' at pos 3, 'e' at pos 4, 'r' at pos 5. 'i' is in the word but not at pos 1. 'n' is in the word but not at pos 2 (since pos 2 was red). So we need a word ending in "ner" (pos 3,4,5) with an 'i' somewhere else (not pos 1) and 'n' not at pos 2. Wait, "inner" has 'n' at 2 and 3. If pos 2 is red and pos 3 is green, then the word has 'n' at 3 but not at 2. And 'i' is present but not at 1. Common words ending in "ner": winner, thinner, dinner, runner, winner. "Winner": w{-}i{-}n{-}n{-}e{-}r. Pos 1 'w' (not 'i'). Pos 2 'i' (yellow in my guess? No, my guess was i{-}n{-}n{-}e{-}r). If target is "winner", guess "inner": i(yellow, should be w), n(red, should be i), n(green, pos 3 is n), e(green, pos 4 is e), r(green, pos 5 is r). This fits perfectly. The word is "winner".} \\
\texttt{guess: winner} \\
            }
        }
    }
    & \\ \\

    \theutterance \stepcounter{utterance}  
    & \multicolumn{4}{p{0.6\linewidth}}{
        \cellcolor[rgb]{0.561,0.737,0.561}{
            \makecell[{{p{\linewidth}}}]{
                \texttt{\tiny{[P1$\rangle$GM]}}
                \texttt{explanation: The feedback for "inner" (i<yellow> n<red> n<green> e<green> r<green>) confirms the word ends in "ner" (positions 3, 4, 5 are correct). The 'n' at position 2 is incorrect, and the 'i' at position 1 is present but misplaced. This suggests the word has an 'i' somewhere other than the start and an 'n' at position 3 but not 2. "Winner" fits this pattern perfectly: w{-}i{-}n{-}n{-}e{-}r. In this guess, 'w' replaces 'i' at the start, 'i' moves to position 2 (which is currently red in my guess, but if the target is "winner", position 2 is 'i', so it should be green? Wait. If target is "winner", guess "inner": pos1 'i' vs 'w' {-}> red? No, 'i' is in the word. If 'i' is in the word but not at pos 1, it gets yellow. Pos 2 'n' vs 'i' {-}> red. Pos 3 'n' vs 'n' {-}> green. Pos 4 'e' vs 'e' {-}> green. Pos 5 'r' vs 'r' {-}> green. This matches the feedback exactly. Therefore, the target is likely "winner". I will guess "winner" to confirm.} \\
\texttt{guess: winner} \\
            }
        }
    }
    & & \\ \\

    \theutterance \stepcounter{utterance}  
    & & & \multicolumn{2}{p{0.3\linewidth}}{
        \cellcolor[rgb]{0.9,0.9,0.9}{
            \makecell[{{p{\linewidth}}}]{
                \texttt{\tiny{[GM$|$GM]}}
                \texttt{Error: The length of the guessed word is not 5.} \\
            }
        }
    }
    & & \\ \\

    \theutterance \stepcounter{utterance}  
    & & & \multicolumn{2}{p{0.3\linewidth}}{
        \cellcolor[rgb]{0.9,0.9,0.9}{
            \makecell[{{p{\linewidth}}}]{
                \texttt{\tiny{[GM$|$GM]}}
                \texttt{game\_result = ABORT} \\
            }
        }
    }
    & & \\ \\

\end{supertabular}
}

\end{document}
